% ======================================================================
% SECTION 6: HR 9506 - American Physical AI Oncology Leadership Act
% of 2026
% New Federal Authorizing Statute Adapting FDORA Section 3209 and
% 21st Century Cures Act 2016.
%
% This file consolidates instructions previously housed in the prior
% template's Implementation/Metrics (subsection 10.1 cross-bill
% timeline diagram, subsection 10.3 guardrail (iv) NEW physical-AI
% guardrail, plus the NEW Domain 5 American Leadership Index
% pointer), Discussion (11.1.d FDORA Section 3209 extension, 11.2.f
% supplementary author simulation evidence, 11.3.a FDA RTCT framing
% closing), Limitations (12.1.f HR 4506 limitation), Future Work
% (Track A, Track B, Track C tracks plus three concrete
% deliverables), and Conclusions (13.1 sixth list item, 13.2 thesis,
% 13.3 themes, 13.4 implications, 13.5 restated limitations, 13.6
% forward path) into the body of HR 9506. HR 9506 is the project
% brief's "United States Number 1" statute and absorbs the cross-
% cutting paper-level material.
% ======================================================================

[CONSOLIDATED INTRODUCTION TIE-IN (formerly Introduction block 2
item 2.d). FDORA / FDA Modernization Act 2.0, PL 117-328 (2022),
Section 3209: in silico and computer-model alternatives to animal
testing; mandate on diversity action plans; mandate on FDA
decentralized trial guidance. HR 9506 closes the gap that no
federal authorization for sustained American leadership in medical
AI and robotics has been enacted as a standalone statute. Cite
\cite{rank4_fdora_fda_modernization_2022} and
\cite{fda_dct_guidance_2024}.]

[CONSOLIDATED CROSS-BILL IMPLEMENTATION TIMELINE (formerly
Implementation Metrics subsection 10.1, plus one comprehensive
ASCII timeline diagram). Compose a verbatim ASCII timeline that
maps each of the seven bills (HR 9501 through HR 9507) to its
12-month or 18-month or 24-month deadline, drawing from each
bill's Subsection .4 in this paper. Use box-drawing characters
(+, -, |) only and keep every line under 80 columns. The
American Leadership Index publication cycle is anchored to this
timeline.]

[PROCESSING INSTRUCTION (HR 9506 - American Physical AI Oncology
Leadership Act of 2026, populate to 12 to 16 paragraphs structured
into the five legislative-act subsections):

The bill is the project brief's "United States Number 1" statute.
It must explicitly state in its first paragraph that it is a "New
Federal Authorizing Statute Adapting FDORA Section 3209
(PL 117-328) and the 21st Century Cures Act of 2016 (PL 114-255)"
so that readers can trace the chain of prior authority. Cite
\cite{rank4_fdora_fda_modernization_2022} and
\cite{rank1_cures_act_2016}. Maintain the respectful FDA framing
throughout: this is an authorization, not a critique.

Subsection 6.1 - Findings and Declarations: Read in full
- national-platform/README.md (the National Platform overview)
- national-platform/new_paper/main.tex and
  national-platform/new_paper/sections/executive_summary.tex,
  introduction.tex, regulatory_landscape.tex, gov_framework.tex,
  implementation_strategy.tex, financial_analysis.tex,
  site_establishment.tex, and conclusion.tex (read all eight)
- patients/paper/Deep-Research-3/part2_ai_robotics_legislation.md
  (the legislative-vehicle taxonomy with the federal authorizing
  statute ranked first)

Compose 6 to 8 enumerated findings citing
\cite{kawchak_2026_19244918}, \cite{main_repo},
\cite{rank4_fdora_fda_modernization_2022}, \cite{fda_rtct_press_2026},
and \cite{fda_oce_advancing_oncology_dct_2024}. Each finding must
state explicitly that the United States must remain Number 1 in
the world for patient benefit during the medical AI and robotics
revolution. At least one finding must reference the FDA's own
April 2026 RTCT announcement as evidence of the FDA's leadership
posture that this bill formalizes into statute.

Subsection 6.2 - Definitions:
- "Medical AI and robotics revolution" means the period from
  approximately 2024 through 2030 during which Physical AI systems
  and large-context AI agents (e.g., Claude Code Opus 4.7 with 1M
  token context, see \cite{claude_opus_47}) are entering routine
  clinical trial deployment. Reference the four-simulation
  evidence base at
  new-trial/national-24-7-trial/paper/full-paper/sections/results.tex.
- "American leadership" means a documented United States lead in:
  (a) the number of authorized Physical AI oncology trial sites,
  (b) the number of patient-self-selected trial bookings per year,
  (c) the cumulative documented error-rate reduction across CMS-
  billable sites, (d) the number of patients whose data is
  governed under HR 9507 self-custody, and (e) the number of FDA
  RTCT pilot trials with patient-sponsor direct channels under HR
  9505.
- "Federal investment" means dedicated appropriations to FDA, NIH
  (PAR-25-170 and successor programs, see
  \cite{nih_par25_170_dht_endpoints}), AHRQ, and ONC for the
  seven-year period 2026 through 2032.

Subsection 6.3 - American Leadership Authorizations (operative
core):
- Authorization for the FDA RTCT pilot program to expand from the
  initial two trials (TRAVERSE, STREAM-SCLC) to at least 50 trials
  by 2028. Cite \cite{fda_rtct_press_2026}.
- Authorization for FDA Oncology Center of Excellence to publish
  per-state dashboards of Physical AI oncology trial site
  authorization counts, harmonized with the prior site
  documentation framework at
  new-trial/site/05-state-regulations/ca_state_regulations.tex
  and national-platform/new_trial_psl/05_title22_regulations.tex.
- Authorization for federated-learning infrastructure across CMS-
  billable sites, drawing on the federated-learning framework at
  national-platform/federated_learning/main_chunk1_preamble_intro_methods.tex,
  main_chunk2_results_architecture_pillars_examples.tex,
  main_chunk3_results_peerreview_trust_analytics.tex, and
  main_chunk4_discussion_limitations_conclusion.tex (read all
  four chunks).
- Authorization for federated investment in foundation models for
  pathology (Virchow), prediction (SCORPIO, PROGPATH), and trial
  matching (TrialGPT, PRISM, TrialMatchAI), with annual per-model
  reporting on patient-benefit metrics. Cite
  \cite{vorontsov2024virchow}, \cite{jin2024trialgpt},
  \cite{gupta2024prism}, and \cite{abdallah2026trialmatchai}.
- Authorization for a public "American Leadership Index"
  published annually with the five metrics defined in 6.2 above.

Subsection 6.4 - Implementation: 12-month deadline for the first
American Leadership Index publication; 24-month deadline for FDA
RTCT pilot expansion to 25 trials (target 50 by 2028); FDA, NIH,
AHRQ, and ONC roles. Reference
national-platform/new_template/sections/implementation_strategy.tex
for the prior implementation pattern.

Subsection 6.5 - Reporting and Enforcement: Annual American
Leadership Index report; Congressional review under FDORA
authorities; explicit alignment with the FDORA Section 3209 in
silico authorization. Cite \cite{rank4_fdora_fda_modernization_2022}.]

[CONSOLIDATED NEW PHYSICAL-AI GUARDRAIL (formerly Implementation
Metrics subsection 10.3 guardrail (iv), the NEW guardrail
introduced by this paper). Physical-AI executor selections and
procedural modifications must always permit a patient-initiated
reversal to human-only execution within the same calendar day,
drawing on patient-journey/master_journey.py and the audit-trail
manager at sponsor/final_paper/scripts/core_agents/audit_trail_manager.py
for the operational baseline. The guardrail is anchored under HR
9506's annual American Leadership Index reporting authority. Cite
\cite{kawchak_2026_19119939} and \cite{kawchak_2026_19396256}.]

[CONSOLIDATED FUTURE WORK TRACKS (formerly Limitations and Future
Work block 2 - Three Future Tracks). HR 9506 authorizes federal
investment that supports three future tracks for delivering all
seven bills' obligations:

Track A - Single big model performing all bills' obligations: a
future generation of Claude Code Opus 4.7 (or equivalent
foundation-model agent) holds all seven bill texts plus the
AVAILABLE DIRECTORIES content in a single 1M-token context and
emits implementation deliverables (machine-readable eligibility
reports, USL transparency reports, modification matrices, error-
rate publications, real-time aggregator endpoints, American
Leadership Index, self-custody endpoints) in a single inference
loop. This track is the natural extension of the four-simulation
paper at new-trial/national-24-7-trial/paper/full-paper/. Cite
\cite{kawchak_2026_19994945} and \cite{claude_opus_47}.

Track B - Distributed agents per bill, coordinated by a smaller
orchestrator: each of the seven bills has its own dedicated agent
(self-selection agent for HR 9501, robot-choice agent for HR 9502,
modification agent for HR 9503, error-rate agent for HR 9504,
sponsor-channel agent for HR 9505, leadership-index agent for HR
9506, self-custody agent for HR 9507), all coordinated by a
smaller orchestrator that runs on hardware as constrained as the
Core i5-6200U with 4 GB RAM documented at
sponsor/final_paper/168_hours/instructions/core_i5_6200u_4gb/README.md.
This track is the natural extension of Simulation 4 at
sponsor/final_paper/168_hours/. Cite \cite{kawchak_2026_19396256}.

Track C - Patient-side native deployment: a future generation
deploys a single integrated patient-side agent that operates HR
9501 through HR 9507 from the patient's own device, using
patient-controlled credentials and HIPAA-compliant authentication
to interact with provider, payer, and sponsor endpoints. This
track addresses the patient priority thesis directly and is the
natural extension of patients/patient_robot_instructions_fixed.tex
pages 1 through 10 plus the digital-twin infrastructure at
patient-journey/stage_03_digital_twin.py. Cite
\cite{kawchak_2026_18810541} and \cite{kawchak_2026_19119939}.]

[CONSOLIDATED FUTURE WORK DELIVERABLES (formerly Limitations and
Future Work block 2 closing paragraph - three concrete
deliverables). HR 9506 authorizes federal investment supporting
three concrete future deliverables:
(i) a TRIPOD+AI / CREMLS-compliant retrospective validation of
HR 9504's error-rate baseline on a real-patient cohort, drawing
on the SHIELD-RT design pattern;
(ii) a public benchmark of HR 9501's machine-readable eligibility
report performance across the TrialGPT, PRISM, and TrialMatchAI
matching systems on a shared dataset;
(iii) an FDA-aligned RTCT pilot submission that uses Simulation 4
(sponsor/final_paper/168_hours/) as the sponsor-side input and HR
9505's patient-sponsor direct channel as the patient-side input.
Cite \cite{fda_rtct_press_2026},
\cite{kawchak_2026_19994945}, \cite{kawchak_2026_19396256}, and
\cite{kawchak_2026_19119939}.]

[CONSOLIDATED DISCUSSION COMPARISON AND THESIS (formerly Discussion
items 11.1.d FDORA Section 3209 extension and 11.3 closing
paragraph plus the original Conclusions thesis 13.2). HR 9506
codifies American leadership under the FDORA in silico
authorization, extending Section 3209 from animal-testing
alternatives to physical-AI executor selection. The central
thesis remains: cancer patients in U.S. oncology clinical trials
are the priority participants and must hold statutory control over
discovery, scheduling, robot selection, procedural modification,
real-time data feedback, error-rate information, and self-custody
of their health data. Advanced AI and advanced robotics now offer
the patient the opportunity to take that control directly; the
seven bills supply the legal infrastructure that makes the
control enforceable and measurable. Cite
\cite{rank4_fdora_fda_modernization_2022},
\cite{kawchak_2026_19994945}, and \cite{kawchak_2026_19119939}.]

[CONSOLIDATED PERSISTENT THEMES (formerly Conclusions item 13.3).
The persistent themes anchored in HR 9506 are:
- Adaption / Revision framing: every bill explicitly states the
  prior U.S. legislative document it builds from.
- Respectful FDA framing: every bill positions the FDA and other
  governing bodies as enabling rather than gating, while remaining
  opportunistic under the patient's control based primarily on
  their own interests.
- Individual patient + broad adoption framing: every bill names
  per-patient examples (PAT-2026-0042 from patient-journey/) AND
  broad-adoption deadlines (12, 18, 24 months across CMS-billable
  sites and 50 states).
- United States Number 1: every bill closes with the leadership
  framing required to keep the United States Number 1 in the
  world for patient benefit during the medical AI and robotics
  revolution.]

[CONSOLIDATED IMPLICATIONS (formerly Conclusions item 13.4). LLM-
scale repository context (1M tokens) demonstrated in
new-trial/national-24-7-trial/paper/full-paper/, the four LLM
simulations summarized at \cite{kawchak_2026_19994945}, the
168-hour sponsor extension at sponsor/final_paper/168_hours/, the
10-stage single-patient journey at patient-journey/, and the
10-robot-category patient instructions at
patients/patient_robot_instructions_fixed.tex collectively show
that the technical substrate is ready. The seven bills supply the
missing legal substrate. Together they enable the cancer patient
to take control of their disease in ways that prior error-prone
and physically and intellectually limited human-only care could
not support.]

[CONSOLIDATED LIMITATIONS PARAGRAPH (formerly Limitations item
12.1.f, plus the cross-cutting paragraph after the per-bill
limitations). HR 9506 is the most ambitious; it assumes seven-
year sustained federal investment which is subject to
appropriations cycles. The American Leadership Index relies on
year-over-year metric stability across 50 states. The bills are
draft language by an independent author and have not been reviewed
by Congress, FDA, or any state legislature. The bills are not
endorsed by Anthropic, OpenAI, Google, or any technology vendor.
The patient-control framework remains an operational hypothesis;
outcome studies are required to confirm that the bills, taken
together, increase participation and decrease error rates without
unintended consequences. The bills do not address pediatric or
vulnerable-population protections beyond the existing 21 CFR 50
Subpart D framework documented at
national-platform/21cfr50_adapt/02_irb_review_pediatric.tex and at
\cite{kawchak_2026_19040707}. Cite
\cite{rank4_fdora_fda_modernization_2022}.]

[CONSOLIDATED FORWARD PATH (formerly Conclusions item 13.6). The
next Claude Code Opus 4.7 1M Max generation pass will populate
this template into a 70+ page paper with full prose, additional
ASCII diagrams, and additional individual patient and robot
examples drawn from the AVAILABLE DIRECTORIES; the author will
then perform the senior-author white-space and table-column-width
formatting pass to remove any orphans, widows, or large empty
page regions before publication. Cite the future Zenodo
deposition at \cite{kawchak_2026_20045457}, the FDA RTCT
announcement at \cite{fda_rtct_press_2026}, the four-simulation
paper at \cite{kawchak_2026_19994945}, the sponsor paper at
\cite{kawchak_2026_19396256}, the National Platform paper at
\cite{kawchak_2026_19244918}, and the patient instructions paper
at \cite{kawchak_2026_18810541}.]

[CONSOLIDATED CONCLUSIONS RESTATEMENT (formerly Conclusions item
13.1 sixth list item). HR 9506 codifies American leadership under
the FDORA in silico authorization, extending Section 3209 from
animal-testing alternatives to physical-AI executor selection.
Cite \cite{rank4_fdora_fda_modernization_2022} and
\cite{rank1_cures_act_2016}.]

Close the bill (1 paragraph after Subsection 6.5) with two
patient-control framings:

(i) FOR INDIVIDUAL PATIENTS: A cancer patient anywhere in the
    United States can verify, in the published American
    Leadership Index, that their state has at least one
    authorized Physical AI oncology trial site, that the per-task
    error rates have improved year-over-year, and that the FDA
    RTCT pilot covers their cancer type. The patient PAT-2026-0042
    in patient-journey/ illustrates the per-patient outcome
    trajectory.

(ii) FOR BROAD ADOPTION: All 50 states must have at least one
     authorized Physical AI oncology trial site within 36 months;
     the United States retains Number 1 status in the world for
     patient benefit during the medical AI and robotics revolution
     by virtue of the per-state, per-cancer, per-site dashboard
     publication; federal appropriations sustain the FDA, NIH,
     AHRQ, and ONC infrastructure for seven years. Cite
     \cite{kawchak_2026_19244918},
     \cite{kawchak_2026_19994945}, and
     \cite{kawchak_2026_19396256}.
